\section{Repeated learning, difficult geometry, and complete deployment cost}\label{sec:r19-experiment}
The empirical design tests the link in Section~\ref{sec:r19-bridge} rather than using the star's easy scalar solve as evidence for difficult stochastic control. We retain the entire earlier information hierarchy, verified-cell studies, star experiments, and adverse nonlinear-transfer results in the companion and the immutable predecessor documents. The current multistage observations were produced in the R16 research run and are integrated into the manuscript here for the first time. Our fleet inference is a retrospective reanalysis of those observations. The accuracy-targeted cost study and deliberate boundary checks are new follow-up computations, separately identified in their execution records.

\subsection{A multistage, vector, nonsmooth study}
The fixed primitive has 63 nodes on a binary tree with six review levels, two service coordinates per node, six coupled quadratic resources, and four additional cross-service capacity constraints. Outside tiers vary by node and service. Every nonroot subtree retains its payment cap, and each service has its root payment equality. Absolute switching costs include installation and have endogenous signs. The displacement formulation is exactly the quadratic instance following Theorem~\ref{thm:r19-inexact}; it is not a collection of independent stars. Each context changes six resource-reward coordinates and a strictly positive friction coefficient. The complete rational primitive and context grids are recorded in the companion.

We evaluate eight independently seeded training/deployment pairs. Each pair uses 256 training contexts and 256 fresh deployment contexts. No best seed or model is selected. All seven methods use the same exact-solver label population: tanh random-feature scalar critics fitted with value-only or value-and-gradient losses; a direct-price tanh regressor; the corresponding three quadratic response surfaces; and a cubic radial-basis direct-price interpolant. The hidden tanh weights are fixed random features, not a fully trained deep network. Value-only and derivative-supervised critics both deploy through their derivative. Direct-price methods predict the same six-dimensional operational statistic. Thus the comparison does not grant neural methods a different response operator or a more favorable certificate.

The feasible outside protocol is specified and node-dependent, not proved to be the reoptimized best static protocol for this heterogeneous primitive. This study tests accepted implementation relative to that announced comparator. The separate canonical and centered-friction studies retain their reoptimized static comparators. Monetary values in the current study are total normalized tree-objective units; they are not per-review gains from the earlier star normalization. The primitive is synthetic and not calibrated to a service provider. Those distinctions prevent a gain against one known feasible protocol from being advertised as an empirical universal value of adaptation.

The R16 solver amendment increased the iteration ceiling after a failure in the R15 execution, while retaining seeds, methods, loss weights, and tolerances; the full amended study was rerun. We describe this as a frozen, independently deployed internal experiment, not an externally preregistered study. Neither this reanalysis nor the follow-up computations modify its trained models or replace its unfavorable observations.

\subsection{Policy quality and the meaning of derivative supervision}
Table~\ref{tab:r19-quality} reports all seven methods on all 2,048 deployment contexts. The reference objective is bracketed by its primal value and independent dual upper bound. Therefore the reported regret is an upper bound, not an unverified subtraction from a numerical optimizer. The average reference bracket width is $8.11\times10^{-8}$.
\begin{table}[p]
\centering
\caption{Unrefined multistage decisions and conditional frozen-fleet gains}\label{tab:r19-quality}
\setlength{\tabcolsep}{4pt}
\begin{tabular}{lrrrr}
\toprule
Method & Regret upper & Mean gap & Pass (\%) & Fleet lower \\
\midrule
Tanh: value & 0.034374 & 0.193705 & 0.00 & 0.542816 \\
Tanh: value + gradient & 0.003685 & 0.015856 & 0.54 & 0.573506 \\
Tanh: price & 0.001581 & 0.007986 & 3.66 & 0.575609 \\
Quadratic: value & 0.003079 & 0.012446 & 4.59 & 0.574111 \\
Quadratic: value + gradient & 0.002794 & 0.011081 & 9.47 & 0.574396 \\
Quadratic: price & 0.000271 & 0.001238 & 65.48 & 0.576919 \\
RBF: price & 0.000158 & 0.000657 & 86.43 & 0.577032 \\
\bottomrule
\end{tabular}

\par\medskip\begin{minipage}{\textwidth}All 2,048 independent deployment contexts per method are included. Pass means the immediate exact full-objective gap is at most $10^{-3}$. No method passes at $10^{-5}$ before refinement. Fleet bounds apply to the explicitly randomized eight-model deployment, not each model or arbitrary retraining.\end{minipage}
\end{table}

The mean regret upper bound falls from $0.0343743$ for value-only tanh fitting to $0.00368462$ with derivative supervision. Across eight independent training/deployment pairs, the approximate paired Student interval for their mean difference is $[-0.0415827,-0.0197968]$. This comparison resolves the direction of improvement for these two fitted procedures in this distribution; it does not overwrite the statistically unresolved R14 star comparison. The independent unit is a complete training/deployment pair, not an individual test context or timing repetition.

Direct-price tanh fitting improves further, to $0.00158094$. Quadratic direct-price and radial-basis fitting attain $0.000271149$ and $0.000157737$, respectively. The latter two also have much higher immediate certificate pass rates at $10^{-3}$. Thus the derivative statistic is operationally useful while neural parameterization is not necessary in this study. This is a result of the required comparison, not a reason to remove neural methods, their mathematical guarantee, or their adverse observations. It identifies the stronger current implementation of the same accepted-control architecture.

\subsection{Conditional gain for an explicitly deployed frozen fleet}
A simultaneous lower bound for each of 56 individual frozen policies, using only its own 256 contexts and the original analytical gain range of 25, is zero. We retain that result. Pooling 2,048 contexts as though they came from one selected model would be invalid. A different, explicitly implementable policy has a more informative guarantee: for each new independent context choose one of the eight frozen policies uniformly, independently of that context, then execute its audited response and outside gate.

\begin{proposition}[Conditional frozen-fleet guarantee]\label{prop:r19-fleet}
For method $m=1,\ldots,M$, fix $K$ policies before inspecting their deployment cohorts. Conditional on all fitted policies, let $G_{mki}\in[0,B]$ be gains on $n$ independent contexts for policy $k$, with all $Kn$ contexts independent across $k,i$. Contexts may be shared across methods. Put $\overline G_m=(Kn)^{-1}\sum_{k,i}G_{mki}$ and let $\mathcal G_m$ be the expected gain of the uniform frozen fleet. With probability at least $1-\delta$ over these deployment cohorts, simultaneously for all methods,
\begin{equation}\label{eq:r19-fleet}
 \mathcal G_m\geq\max\left\{0,\overline G_m-
 B\sqrt{\frac{\log(M/\delta)}{2Kn}}\right\}.
\end{equation}
This does not assert the bound for each constituent policy or for a population of newly trained policies.
\end{proposition}
\begin{proof}
Conditional on the policies, the $Kn$ summands for each method are independent bounded variables; identical distributions across policies are unnecessary. Their expected average is exactly the uniform fleet's expected gain. Hoeffding's lower-tail bound and a union bound over $M$ methods prove the simultaneous inequality. The comparator gate gives nonnegative gain, justifying truncation of the lower bound at zero. Independence across methods is not used.
\end{proof}
For $K=8$, $n=256$, $M=7$, $B=25$, and $\delta=0.05$, the penalty is $0.868352$. The resulting seven lower bounds, included in Table~\ref{tab:r19-quality}, are positive. They certify this finite, frozen randomization rule under the stated context distribution. The eight-pair Student comparison addresses training/deployment variability approximately; the fleet inequality instead gives a finite-sample, conditional guarantee. Neither substitutes for a guarantee over arbitrary retraining, misspecification, or distribution shift.

\subsection{Geometry rather than successful-region averages}
The earlier R16 structural suite had no binding extra capacity among its 64 tested instances, despite some tier boundaries. We do not relabel those instances as a capacity stress test. The new 64-instance suite deliberately includes heterogeneous controls, root cross-service capacities with positive incentives, tier-boundary incentives with slack extra capacities, and high-friction switching kinks. Each stratum uses four primitive seeds and four contexts per seed on 15- or 31-node, two-service trees with two or four resource factors. All 16 cases in each targeted stratum must actually exhibit the specified geometry or the test fails. No failing instance may be dropped.
\begin{table}[p]
\centering
\caption{Deliberate multistage geometry checks}\label{tab:r19-geometry}
\setlength{\tabcolsep}{4pt}
\begin{tabular}{lrrrrr}
\toprule
Stratum & Cases & Capacity & Tier boundary & Kinks & Gradient error \\
\midrule
heterogeneous & 16 & 0 & 0 & 16 & 1.04e-06 \\
binding capacity & 16 & 16 & 0 & 16 & 5.10e-08 \\
tier boundary & 16 & 0 & 16 & 16 & 1.48e-11 \\
switching kinks & 16 & 0 & 0 & 16 & 7.47e-05 \\
\bottomrule
\end{tabular}

\par\medskip\begin{minipage}{\textwidth}Capacity, tier boundary, and kinks count instances with at least one such event; diagnostics use $10^{-5}$. Each of 128 implemented reference/proposal certificates is independently replayed. The gradient column is the maximum finite-difference error among the four checked contexts in that stratum.\end{minipage}
\end{table}

Every reference and repaired response is independently replayed in rational arithmetic. Theorem~\ref{thm:r19-inexact} is checked with its certified response error rather than by treating the numerical subproblem solution as exact. An independent SLSQP solve starts from the outside protocol on a capacity-binding instance. Finite-difference checks perturb only the resource-reward coordinate, keeping acceptance, switching, and all other primitives fixed. Exact rational scalar checks additionally exercise a quartic smooth resource coupling with nonsmooth switching. These are theorem tests, not new trained-policy observations or evidence of calibrated operations performance.

The retained star reassessment also includes all original, critical-friction, capacity-activation, and primitive-perturbation strata. Its capacity-activation stratum binds the root constraint in every case. All methods then require capacity response adjustment and yield the same certified solution to numerical precision. That is evidence for the response solver, not successful neural extrapolation. The companion reports stratification by friction distance, resource binding, leaf activation and saturation, nearest breakpoint, and optimum gain, without selecting only favorable cells.

\subsection{Accuracy-targeted total cost and amortization}
The original R16 timings charge prediction, the base response, rational repair and audit, actual full-objective refinement, and a second audit. They use a $10^{-9}$ numerical tolerance even for the coarse deployment target. That design may over-solve the classical comparator. We preserve those measurements and execute the previously recorded, but unimplemented, accuracy-targeted comparison instead of treating the earlier apparent speed advantage as decisive.

In the follow-up, both classical solvers and all surrogate response solvers first use numerical tolerance $10^{-5}$. Every method must pass the same independent full-objective certificate at the requested tolerance; otherwise it performs a full $10^{-9}$ refinement and another audit. The cached classical baselines use cold and previous-context starts. All seven frozen surrogate models are included. We use the first 16 contexts of every deployment cohort, two repetitions, and rotated method order. The total comparison contains 4,608 completed pipelines, with 256 records per method and requested tolerance. Timing observations and model-selection units are kept separate.
\begin{table}[p]
\centering
\caption{Accuracy-targeted complete-pipeline cost}\label{tab:r19-cost}
\setlength{\tabcolsep}{4pt}
\begin{tabular}{lrrrr}
\toprule
Method & \shortstack{Mean ms\\$10^{-3}$} & \shortstack{Mean ms\\$10^{-5}$} & \shortstack{Fall back\\(\%)} & \shortstack{Fall back\\(\%)} \\
\midrule
Classical: cold & 14.029 & 26.180 & 0.00 & 42.19 \\
Classical: previous & 13.453 & 26.822 & 0.00 & 43.36 \\
Tanh: value & 43.712 & 43.532 & 100.00 & 100.00 \\
Tanh: value + gradient & 42.169 & 42.721 & 99.22 & 100.00 \\
Tanh: price & 41.519 & 41.594 & 98.44 & 100.00 \\
Quadratic: value & 40.264 & 41.201 & 95.31 & 100.00 \\
Quadratic: value + gradient & 39.654 & 42.080 & 91.41 & 100.00 \\
Quadratic: price & 19.790 & 39.858 & 44.14 & 100.00 \\
RBF: price & 9.563 & 37.808 & 12.89 & 100.00 \\
\bottomrule
\end{tabular}

\par\medskip\begin{minipage}{\textwidth}The last two columns correspond to $10^{-3}$ and $10^{-5}$, respectively. All methods start with numerical tolerance $10^{-5}$, then refine and re-audit when the independent full-objective gap fails the requested target. Each mean uses 256 records. No prediction, response, repair, audit, or actual refinement cost is omitted.\end{minipage}
\end{table}

At the coarse target, the lower-mean classical baseline costs 13.453 ms, compared with 9.563 ms for radial-basis direct-price deployment. Its observed full-pipeline saving has a conservative offline break-even count of 2,546 queries. At the strict target, the corresponding classical cost is 26.180 ms and the radial-basis cost is 37.808 ms. These are separately measured follow-up results, not the original fixed-solver-tolerance timings.

The table reports mean complete cost, not the sum of separately reported medians. The companion supplies medians, 95th percentiles, fallback rates, and paired cohort intervals. Offline cost is measured on the same host by replaying all 256 training-label computations and all seven fits for every cohort, while leaving the deployed frozen models unchanged. We conservatively charge that entire shared offline cost to each surrogate. A finite break-even count is reported only when its observed mean online saving is positive. It is a benchmark-specific accounting calculation, not a future runtime guarantee. The data do not establish universal neural superiority, and no failed cost comparison is hidden behind a warm-start or an unaudited prediction time.
